%%=============================================================================
%% Samenvatting
%%=============================================================================

% TODO: De "abstract" of samenvatting is een kernachtige (~ 1 blz. voor een
% thesis) synthese van het document.
%
% Een goede abstract biedt een kernachtig antwoord op volgende vragen:
%
% 1. Waarover gaat de bachelorproef?
% 2. Waarom heb je er over geschreven?
% 3. Hoe heb je het onderzoek uitgevoerd?
% 4. Wat waren de resultaten? Wat blijkt uit je onderzoek?
% 5. Wat betekenen je resultaten? Wat is de relevantie voor het werkveld?
%
% Daarom bestaat een abstract uit volgende componenten:
%
% - inleiding + kaderen thema
% - probleemstelling
% - (centrale) onderzoeksvraag
% - onderzoeksdoelstelling
% - methodologie
% - resultaten (beperk tot de belangrijkste, relevant voor de onderzoeksvraag)
% - conclusies, aanbevelingen, beperkingen
%
% LET OP! Een samenvatting is GEEN voorwoord!

%%---------- Nederlandse samenvatting -----------------------------------------
%
% TODO: Als je je bachelorproef in het Engels schrijft, moet je eerst een
% Nederlandse samenvatting invoegen. Haal daarvoor onderstaande code uit
% commentaar.
% Wie zijn bachelorproef in het Nederlands schrijft, kan dit negeren, de inhoud
% wordt niet in het document ingevoegd.

%%---------- Samenvatting -----------------------------------------------------
% De samenvatting in de hoofdtaal van het document

\chapter*{\IfLanguageName{dutch}{Samenvatting}{Abstract}}

Mainframes spelen nog steeds een belangrijke rol in sectoren zoals banken, verzekeringen en luchtvaart. Veel van de software op deze systemen is geschreven in High Level Assembler (HLASM). Deze programma’s zijn vaak moeilijk te begrijpen, zeker voor beginners. Daarnaast is er een tekort aan nieuwe ontwikkelaars, wat het doorgeven van kennis moeilijk maakt.

Deze bachelorproef onderzoekt hoe visualisatie kan helpen bij het begrijpen van de control-flow van HLASM-programma’s. De centrale onderzoeksvraag luidt: “Hoe effectief versnellen statische analyse en Mermaid-diagrammen het leerproces van HLASM bij beginners?”. Het doel van het onderzoek is het ontwikkelen en evalueren van een proof-of-concept (PoC) die HLASM-code omzet naar een visuele voorstelling.

Het onderzoek werd uitgevoerd in verschillende stappen. Eerst werd een literatuurstudie gedaan naar mainframes, HLASM en bestaande visualisatietechnieken. Daarna werd HLASM-code geanalyseerd om te bepalen welke onderdelen belangrijk zijn voor de control-flow, zoals labels en sprongen. Op basis daarvan werd een proof-of-concept ontwikkeld als webapplicatie. Deze tool zet HLASM-code via statische analyse om naar een flowchart in Mermaid.js. Tot slot werd de tool getest met vier beginnende gebruikers.

Uit de resultaten blijkt dat de visualisatie vooral helpt om sneller een overzicht te krijgen van de structuur van een programma. De deelnemers konden gemakkelijker zien hoe verschillende delen van de code met elkaar verbonden zijn en waar vertakkingen plaatsvinden.

Tegelijk werd duidelijk dat de tool het begrijpen van de code niet volledig vervangt. Om de werking van het programma correct te begrijpen, blijft het nodig om de code zelf te lezen en te analyseren. De visualisatie werd dus vooral gezien als een hulpmiddel dat het proces ondersteunt.

Daarnaast bleek dat de meerwaarde van de tool groter is bij grotere programma’s. Bij kleinere codefragmenten is de structuur meestal eenvoudig genoeg om manueel te volgen, waardoor de visualisatie minder toevoegt.

De resultaten tonen aan dat statische analyse in combinatie met een visuele voorstelling een nuttig hulpmiddel kan zijn voor beginnende ontwikkelaars. De ontwikkelde PoC heeft echter beperkingen: slechts een deel van de HLASM-instructies wordt ondersteund en complexe situaties zoals macro’s en indirecte sprongen worden niet meegenomen.

\newpage

Samenvattend kan gesteld worden dat een eenvoudige visualisatietool helpt om de structuur van HLASM-programma’s sneller te begrijpen. De tool ondersteunt het leerproces, maar vervangt het niet, en is vooral nuttig bij grotere programma’s.